\begin{table*}[!t]
    \centering
    \caption{Summary of trigger value-selection strategies.}
    \label{tab:value_selectors}
    \begin{tabularx}{\textwidth}{>{\raggedright\arraybackslash}p{2.2cm}
    >{\raggedright\arraybackslash}p{3.2cm}
    >{\raggedright\arraybackslash}X
    >{\raggedright\arraybackslash}p{3.2cm}}
        \toprule
        \textbf{Source} & \textbf{Selector} & \textbf{Main principle} & \textbf{Category} \\
        \midrule
        Severi-derived & MinPopulation & Selects rare observed values. & Frequency / rarity \\
        Severi-derived & CountAbsSHAP & Balances value frequency and low absolute SHAP contribution. & Frequency + SHAP magnitude \\
        Severi-derived & CombinedSHAP & Greedily selects benign-oriented co-occurring feature-value pairs. & Signed SHAP + co-occurrence \\
        This work & FreqBoundedSigned & Selects benign-oriented values within a frequency range. & Frequency + signed SHAP \\
        This work & LowSignedSHAP & Selects values with a benign-oriented signed contribution. & Signed SHAP \\
        This work & CorrCountAbsSHAP & Considers low absolute SHAP values with co-occurrence support. & SHAP magnitude + co-occurrence \\
        This work & BenignPrototype & Selects values from representative benign patterns. & Prototype / co-occurrence \\
        This work & QuantileDistribution & Selects values from fixed empirical distribution positions. & Distribution baseline \\
        \bottomrule
    \end{tabularx}
\end{table*}